\chapter{Connectivity}\label{ch:connectivity}

Chapter~\ref{ch:graphs} asked whether two vertices are connected. This chapter asks how \emph{robustly} they are connected, and
what the connections look like at a finer scale. In a directed graph, which vertices can reach each other in both directions?
In an undirected graph, which single edge or vertex is so important that its failure disconnects the network? Can a route use
every road exactly once? And what if the graph changes, with roads opening and closing over time? All of these questions are
answered by variations of one idea, the \emph{low-link} values of a depth-first search, plus a handful of beautiful reductions:
boolean formulas to strongly connected components (2-SAT), tours to degrees (Euler), and deletions to rollbacks (offline dynamic
connectivity).

\begin{chapterintro}{\plat\ \ioi}{Chapters~\ref{ch:graphs} and~\ref{ch:dsu}, in particular DFS trees and edge
classification (Section~\ref{sec:dfs}).}
  \item find strongly connected components with Tarjan's algorithm, and do DP on the condensation;
  \item solve 2-SAT through the implication graph;
  \item find bridges and articulation points with low-link values, and build bridge trees and block-cut trees;
  \item decide whether Eulerian paths and circuits exist, and construct them with Hierholzer's algorithm;
  \item handle edge deletions offline with a segment tree on time and a DSU with rollback;
  \item use Euler's formula for planar graphs, and reason about critical vertices and dominators.
\end{chapterintro}

% ======================================================================================================================
\section[Strongly connected components]{Strongly connected components \hfill \plat}\label{sec:scc}

In a directed graph, two vertices are \defn{strongly connected} if each can reach the other. This is an equivalence relation,
and its classes are the \defn{strongly connected components} (SCCs). Contracting every SCC to a single vertex gives the
\defn{condensation}, which is always a DAG: a cycle through several components would merge them into one
(Figure~\ref{fig:scc}). So every directed graph is a DAG of strongly connected pieces, and many questions split into a question
inside each piece (where everything is mutually reachable) and a DP over the DAG.

\begin{figure}[htbp]
\centering
\begin{tikzpicture}
  \begin{scope}
    \fill[fgsky!18, rounded corners=8pt] (-0.5, -0.25) rectangle (1.8, 2.05);
    \fill[fgorange!15, rounded corners=8pt] (2.7, -0.25) rectangle (3.75, 2.05);
    \fill[fgpink!18, rounded corners=8pt] (4.55, 0.45) rectangle (5.45, 1.55);
    \node[vtx, cool] (a) at (0, 1) {0}; \node[vtx, cool] (b) at (1.3, 1.6) {1}; \node[vtx, cool] (c) at (1.3, 0.3) {2};
    \node[vtx, warn] (d) at (3.2, 1.6) {3}; \node[vtx, warn] (e) at (3.2, 0.3) {4};
    \node[vtx, pur] (f) at (5, 1) {5};
    \draw[ded] (a) -- (b); \draw[ded] (b) -- (c); \draw[ded] (c) -- (a); \draw[ded] (b) -- (d);
    \draw[ded] (d) to[bend left=20] (e); \draw[ded] (e) to[bend left=20] (d); \draw[ded] (c) -- (e); \draw[ded] (e) -- (f);
  \end{scope}
  \draw[->, line width=1.2pt, draw=fgblue] (5.9, 1) -- (7.0, 1);
  \begin{scope}[xshift=7.8cm]
    \node[vtx, cool, minimum size=11mm] (A) at (0, 1) {\scriptsize 0, 1, 2};
    \node[vtx, warn, minimum size=11mm] (B) at (2, 1) {\scriptsize 3, 4};
    \node[vtx, pur, minimum size=11mm] (C) at (4, 1) {\scriptsize 5};
    \draw[ded] (A) -- (B); \draw[ded] (B) -- (C);
    \node[lbl] at (2, -0.1) {the condensation is a DAG};
  \end{scope}
\end{tikzpicture}
\caption{A directed graph with three strongly connected components (shaded), and its condensation.}
\label{fig:scc}
\end{figure}

\paragraph{Tarjan's algorithm.} One DFS finds all SCCs. Every vertex gets a discovery time $\mathit{tin}(v)$ and a
\defn{low-link} value $\mathit{low}(v)$: the smallest discovery time reachable from $v$'s DFS subtree by following tree edges and
then \emph{one} edge to a vertex that is still on a stack of ``open'' vertices. Vertices are pushed when they are discovered. When
DFS finishes a vertex $u$ with $\mathit{low}(u) = \mathit{tin}(u)$, nothing in its subtree can reach an earlier open vertex, so $u$
is the first vertex of its component that DFS discovered, and the component is exactly the vertices above $u$ on the stack: pop
them (Figure~\ref{fig:tarjan}). The whole algorithm is one DFS, $\Oh{n + m}$.

\begin{figure}[htbp]
\centering
\begin{tikzpicture}
  \fill[fgsky!18, rounded corners=8pt] (-0.55, -1.95) rectangle (2.15, 0.55);
  \fill[fgorange!15, rounded corners=8pt] (3.05, -1.95) rectangle (5.75, 0.55);
  \fill[fgpink!18, rounded corners=8pt] (6.65, -1.95) rectangle (7.55, -0.95);
  \node[vtx] (v1) at (0, 0) {a}; \node[vtx] (v2) at (1.6, 0) {b}; \node[vtx] (v3) at (0.8, -1.4) {c};
  \node[vtx] (v4) at (3.6, 0) {d}; \node[vtx] (v5) at (5.2, 0) {e}; \node[vtx] (v6) at (4.4, -1.4) {f};
  \node[vtx] (v7) at (7.1, -1.4) {g};
  \draw[bded] (v1) -- (v2); \draw[bded] (v2) -- (v3); \draw[ded, draw=fgverm, dashed] (v3) -- (v1);
  \draw[bded] (v3) -- (v4); \draw[bded] (v4) -- (v5); \draw[bded] (v5) -- (v6);
  \draw[ded, draw=fgverm, dashed] (v6) -- (v4); \draw[bded] (v6) -- (v7);
  \foreach \vv/\tl/\pos/\an in {v1/{1/1}/left/west, v2/{2/1}/right/east, v3/{3/1}/left/west,
      v4/{4/4}/above left/north west, v5/{5/4}/right/east, v6/{6/4}/left/west, v7/{7/7}/below/south}
    \node[font=\sffamily\scriptsize, text=fgink!70, \pos=0pt] at (\vv.\an) {\tl};
  \node[font=\sffamily\scriptsize\bfseries, text=fgblue] at (0.8, 0.8) {SCC 3};
  \node[font=\sffamily\scriptsize\bfseries, text=fgorange!80!black] at (4.4, 0.8) {SCC 2};
  \node[font=\sffamily\scriptsize\bfseries, text=fgpink!80!black] at (7.1, -0.7) {SCC 1};
\end{tikzpicture}
\caption{Tarjan's algorithm on a graph with three SCCs; labels are $\mathit{tin}/\mathit{low}$, blue edges are tree edges and red
dashed edges go back to open vertices. Vertex $g$ finishes first with $\mathit{low} = \mathit{tin}$ and forms SCC 1; then
vertex $d$ (popping $f$, $e$, $d$), then vertex $a$. Components come out in reverse topological order of the condensation.}
\label{fig:tarjan}
\end{figure}

\begin{problem}{Planets and Kingdoms (CSES)}
Divide the planets of a directed teleporter network into kingdoms, so that two planets are in the same kingdom exactly when
each can reach the other.
\end{problem}

\begin{code}
void solve() {
    int n, m;
    cin >> n >> m;
    vvi adj(n + 1);
    rep(i, m) {
        int a, b;
        cin >> a >> b;
        adj[a].push_back(b);
    }
    // Tarjan: tin = discovery time, low = smallest tin reachable from the
    // subtree through one edge to a vertex that is still on the stack
    vi tin(n + 1, 0), low(n + 1, 0), comp(n + 1, 0), st;
    vector<bool> on(n + 1, false);
    int timer = 0, k = 0;
    auto dfs = [&](auto&& self, int u) -> void {
        tin[u] = low[u] = ++timer;
        st.push_back(u), on[u] = true;
        for (int v : adj[u]) {
            if (!tin[v]) self(self, v), low[u] = min(low[u], low[v]);
            else if (on[v]) low[u] = min(low[u], tin[v]);
        }
        if (low[u] == tin[u]) {  // u is the first vertex of its SCC: pop it
            k++;
            int v;
            do {
                v = st.back();
                st.pop_back();
                on[v] = false, comp[v] = k;
            } while (v != u);
        }
    };
    rep(v, 1, n + 1) if (!tin[v]) dfs(dfs, v);
    print(k);
    print(vi(comp.begin() + 1, comp.end()));
}
\end{code}

\usetemplate{SCC (Tarjan)}{Ctrl+Alt+T, C}{cpscc}

\paragraph{Why components come out in reverse topological order.} If there is an edge from component $A$ to component $B$, then
$B$ is completed before $A$: either DFS entered $B$ first (and $B$, unable to reach $A$, finished before $A$ was even discovered)
or it entered $A$ first and reached $B$ inside $A$'s subtree. So Tarjan's numbers decrease along edges of the condensation, and
processing components in \emph{increasing} number visits every component after all its successors. \defn{Kosaraju's algorithm}
is the alternative: one DFS for finishing times, then DFS on the reversed graph in decreasing finishing time, each tree being an
SCC. It is easier to prove but does twice the work.

\begin{problem}{Flight Routes Check (CSES)}
Decide whether every city can reach every other by one-way flights; if not, print two cities $a, b$ with no route from $a$ to
$b$.
\solution No SCC algorithm is needed: the graph is strongly connected exactly when vertex 1 reaches every vertex and every vertex
reaches 1 (search the reversed graph from 1). A vertex that fails either test gives the pair.
\end{problem}

\begin{code}
void solve() {
    int n, m;
    cin >> n >> m;
    vvi adj(n + 1), radj(n + 1);
    rep(i, m) {
        int a, b;
        cin >> a >> b;
        adj[a].push_back(b);
        radj[b].push_back(a);
    }
    // strongly connected iff 1 reaches everyone and everyone reaches 1
    auto reach = [&](const vvi& g) {
        vector<bool> seen(n + 1, false);
        vi st{1};
        seen[1] = true;
        while (!st.empty()) {
            int u = st.back();
            st.pop_back();
            for (int v : g[u])
                if (!seen[v]) seen[v] = true, st.push_back(v);
        }
        return seen;
    };
    vector<bool> fwd = reach(adj), bwd = reach(radj);
    rep(v, 1, n + 1) {
        if (!fwd[v]) {  // no route 1 -> v
            print("NO");
            print(1, v);
            return;
        }
        if (!bwd[v]) {  // no route v -> 1
            print("NO");
            print(v, 1);
            return;
        }
    }
    print("YES");
}
\end{code}

\begin{problem}{Coin Collector (CSES)}
Each room contains coins, and tunnels are one-way. Starting anywhere and ending anywhere, collect as many coins as possible (a
room's coins are collected once).
\solution Inside an SCC, a walk can visit every room and return, so entering a component means collecting all of it. Condense,
give each component the sum of its coins, and find the heaviest path in the condensation. Tarjan's numbering is already a reverse
topological order, so one pass in increasing component number computes the best path starting at each component.
\end{problem}

\begin{code}
void solve() {
    int n, m;
    cin >> n >> m;
    vi coins(n + 1);
    rep(v, 1, n + 1) cin >> coins[v];
    vvi adj(n + 1);
    rep(i, m) {
        int a, b;
        cin >> a >> b;
        adj[a].push_back(b);
    }
    vi tin(n + 1, 0), low(n + 1, 0), comp(n + 1, 0), st;
    vector<bool> on(n + 1, false);
    int timer = 0, k = 0;
    auto dfs = [&](auto&& self, int u) -> void {
        tin[u] = low[u] = ++timer;
        st.push_back(u), on[u] = true;
        for (int v : adj[u]) {
            if (!tin[v]) self(self, v), low[u] = min(low[u], low[v]);
            else if (on[v]) low[u] = min(low[u], tin[v]);
        }
        if (low[u] == tin[u]) {
            k++;
            int v;
            do {
                v = st.back();
                st.pop_back();
                on[v] = false, comp[v] = k;
            } while (v != u);
        }
    };
    rep(v, 1, n + 1) if (!tin[v]) dfs(dfs, v);
    // Tarjan numbers SCCs in reverse topological order: an edge between
    // components goes from a larger number to a smaller one
    vi total(k + 1, 0), best(k + 1, 0);
    rep(v, 1, n + 1) total[comp[v]] += coins[v];
    vvi out(k + 1);
    rep(u, 1, n + 1) for (int v : adj[u])
        if (comp[u] != comp[v]) out[comp[u]].push_back(comp[v]);
    ll ans = 0;
    rep(c, 1, k + 1) {  // successors (smaller numbers) are already done
        for (int d : out[c]) chmax(best[c], best[d]);
        best[c] += total[c];
        chmax(ans, best[c]);
    }
    print(ans);
}
\end{code}

\begin{recognize}
``Mutually reachable'', ``can every vertex reach every other?'', cyclic dependencies, and ``maximise what you collect when you may
revisit vertices'' (condense, then DP on the DAG). Also the minimum number of edges to add to make a graph strongly connected:
$\max(\#\text{sources}, \#\text{sinks})$ of the condensation, unless it is already a single component.
\end{recognize}

% ======================================================================================================================
\section[2-SAT]{2-SAT \hfill \plat}\label{sec:2sat}

A \defn{2-SAT} instance has boolean variables and clauses of the form $(a \lor b)$, where $a$ and $b$ are variables or negated
variables (\emph{literals}). General satisfiability is NP-hard, but with two literals per clause it is solvable in linear time.
The clause $(a \lor b)$ is equivalent to the two implications $\lnot a \Rightarrow b$ and $\lnot b \Rightarrow a$. Build the
\defn{implication graph} on the $2n$ literals with these edges (Figure~\ref{fig:2sat}). Then:

\begin{itemize}
  \item if some variable $x$ and its negation $\lnot x$ lie in the same SCC, then $x \Rightarrow \lnot x$ and $\lnot x \Rightarrow
        x$, and no assignment works;
  \item otherwise, set $x$ true exactly when the SCC of $x$ comes \emph{after} the SCC of $\lnot x$ in topological order (with
        Tarjan's numbering: when $\mathit{comp}(x) < \mathit{comp}(\lnot x)$).
\end{itemize}

\noindent The second rule works because the implication graph is symmetric (an edge $u \to v$ comes with $\lnot v \to \lnot u$):
choosing, for each variable, the literal later in topological order never creates an implication from a true literal to a false
one.

\begin{figure}[htbp]
\centering
\begin{tikzpicture}
  \fill[fggreen!15, rounded corners=10pt] (2.3, -0.6) rectangle (6.1, 2.2);
  \fill[fgverm!10, rounded corners=10pt] (-1.9, -0.6) rectangle (1.9, 2.2);
  \node[vtx, minimum size=8mm] (na) at (-1.2, 1.5) {$\lnot a$};
  \node[vtx, minimum size=8mm] (b) at (1.2, 1.5) {$b$};
  \node[vtx, minimum size=8mm] (nc) at (0, 0) {$\lnot c$};
  \node[vtx, minimum size=8mm] (a) at (3.0, 1.5) {$a$};
  \node[vtx, minimum size=8mm] (c) at (5.4, 1.5) {$c$};
  \node[vtx, minimum size=8mm] (nb) at (4.2, 0) {$\lnot b$};
  \draw[ded] (na) -- (b); \draw[ded] (b) -- (nc); \draw[ded] (nc) -- (na);
  \draw[ded] (a) -- (c); \draw[ded] (c) -- (nb); \draw[ded] (nb) -- (a);
  \draw[hded] (na) to[bend left=35] (a);
  \node[font=\sffamily\scriptsize\bfseries, text=fgverm] at (0, -0.95) {false side};
  \node[font=\sffamily\scriptsize\bfseries, text=fggreen!50!black] at (4.2, -0.95) {true side};
  \node[notel, anchor=west, text width=5.0cm] at (6.5, 0.8) {Clauses $(a \lor b)$, $(\lnot a \lor c)$, $(\lnot b \lor \lnot c)$,
    $(a \lor a)$. Each clause gives two edges; $(a \lor a)$ gives $\lnot a \to a$ (red). The SCC of $a$ comes later in
    topological order than that of $\lnot a$, so $a$, $c$ and $\lnot b$ are true.};
\end{tikzpicture}
\caption{An implication graph with two SCCs that are mirror images of each other. The edge between them fixes the order, and
the later component is set true: $a = \text{true}$, $b = \text{false}$, $c = \text{true}$.}
\label{fig:2sat}
\end{figure}

\begin{problem}{Giant Pizza (CSES)}
Each of $n$ family members has two wishes of the form ``topping $x$ on'' or ``topping $x$ off'', and wants at least one of them
fulfilled. Choose the $m$ toppings to satisfy everyone, or report that it is impossible.
\end{problem}

\begin{code}
void solve() {
    int n, m;  // n wishes, m toppings
    cin >> n >> m;
    // literal 2x: topping x is on, 2x + 1: it is off
    int N = 2 * m;
    vvi adj(N);
    auto lit = [](char sign, int x) { return 2 * (x - 1) + (sign == '-'); };
    rep(i, n) {
        char s1, s2;
        int x1, x2;
        cin >> s1 >> x1 >> s2 >> x2;
        int a = lit(s1, x1), b = lit(s2, x2);
        adj[a ^ 1].push_back(b);  // (a or b): not a implies b,
        adj[b ^ 1].push_back(a);  // and not b implies a
    }
    vi tin(N, 0), low(N, 0), comp(N, 0), st;
    vector<bool> on(N, false);
    int timer = 0, k = 0;
    auto dfs = [&](auto&& self, int u) -> void {
        tin[u] = low[u] = ++timer;
        st.push_back(u), on[u] = true;
        for (int v : adj[u]) {
            if (!tin[v]) self(self, v), low[u] = min(low[u], low[v]);
            else if (on[v]) low[u] = min(low[u], tin[v]);
        }
        if (low[u] == tin[u]) {
            k++;
            int v;
            do {
                v = st.back();
                st.pop_back();
                on[v] = false, comp[v] = k;
            } while (v != u);
        }
    };
    rep(v, N) if (!tin[v]) dfs(dfs, v);
    string ans(m, '?');
    rep(x, m) {
        if (comp[2 * x] == comp[2 * x + 1]) {  // x on implies x off and back
            print("IMPOSSIBLE");
            return;
        }
        // choose the literal whose SCC is later in topological order,
        // which is the smaller Tarjan number
        ans[x] = comp[2 * x] < comp[2 * x + 1] ? '+' : '-';
    }
    rep(x, m) cout << ans[x] << (x + 1 < m ? ' ' : '\n');
}
\end{code}

\usetemplate{2-SAT}{Ctrl+Alt+Y, 2}{cp2sat}

\begin{recognize}
Every object has exactly two options (on/off, left/right, team A/B, top/bottom), and every constraint talks about at most two
objects: ``not both'' is $(\lnot a \lor \lnot b)$, ``at least one'' is $(a \lor b)$, ``exactly one'' is both, ``if $a$ then
$b$'' is $(\lnot a \lor b)$, and ``$a$ must be true'' is $(a \lor a)$. Constraints on three objects at once make the problem
NP-hard, so look for a way to express every constraint pairwise. ``At most one of these $k$'' is expressible with $\Oh{k}$
auxiliary prefix variables instead of $\Oh{k^2}$ clauses.
\end{recognize}

% ======================================================================================================================
\section[Bridges and articulation points]{Bridges and articulation points \hfill \plat}\label{sec:bridges}

In an undirected graph, a \defn{bridge} is an edge whose removal increases the number of connected components, and an
\defn{articulation point} (cut vertex) is a vertex whose removal does. Both are found with one DFS and low-link values: here
$\mathit{low}(v)$ is the smallest discovery time reachable from $v$'s subtree using tree edges downwards and then at most one back
edge (Figure~\ref{fig:bridges}). Recall that an undirected DFS has only tree and back edges, and back edges go to ancestors.

\begin{itemize}
  \item A tree edge $u$--$v$ ($v$ the child) is a \textbf{bridge} if and only if $\mathit{low}(v) > \mathit{tin}(u)$: nothing in
        $v$'s subtree has a back edge to $u$ or above, so removing the edge cuts the subtree off. Back edges are never bridges;
        they lie on a cycle.
  \item A non-root vertex $u$ is an \textbf{articulation point} if and only if some child $v$ has $\mathit{low}(v) \ge
        \mathit{tin}(u)$: $v$'s subtree cannot get above $u$ without $u$. The root is one exactly when it has at least two DFS
        children.
\end{itemize}

\begin{figure}[htbp]
\centering
\begin{tikzpicture}
  \node[vtx] (a) at (0, 0) {a}; \node[vtx] (b) at (1.3, 0.9) {b}; \node[vtx] (c) at (1.3, -0.9) {c};
  \node[vtx, hot] (d) at (2.8, 0) {d}; \node[vtx, hot] (e) at (4.4, 0) {e};
  \node[vtx] (f) at (5.8, 0.9) {f}; \node[vtx] (g) at (5.8, -0.9) {g}; \node[vtx] (h) at (7.3, 0) {h};
  \draw[bed] (a) -- (b); \draw[bed] (b) -- (d); \draw[bed] (d) -- (c); \draw[xed] (c) -- (a);
  \draw[hed] (d) -- (e);
  \draw[bed] (e) -- (f); \draw[bed] (f) -- (h); \draw[bed] (h) -- (g); \draw[xed] (g) -- (e); \draw[xed] (f) -- (g);
  \foreach \vv/\tl/\pos/\an in {a/{1/1}/left/west, b/{2/1}/above/north, c/{4/1}/below/south,
      d/{3/1}/below right/south east, e/{5/5}/below left/south west, f/{6/5}/above/north, g/{8/5}/below/south,
      h/{7/5}/right/east}
    \node[font=\sffamily\scriptsize, text=fgink!70, \pos=0pt] at (\vv.\an) {\tl};
  \node[notel, anchor=west, text width=4.1cm] at (8.7, 0) {Labels $\mathit{tin}/\mathit{low}$ from a DFS started at $a$ (blue
    tree edges, grey dashed back edges). $\mathit{low}(e) = 5 > \mathit{tin}(d) = 3$, so $d$--$e$ is a bridge; $d$ and $e$
    are articulation points.};
\end{tikzpicture}
\caption{Bridges and articulation points from low-link values: the red edge is the only bridge, and the red vertices are the
articulation points.}
\label{fig:bridges}
\end{figure}

\begin{problem}{Necessary Roads (CSES)}
List every road whose removal would disconnect some pair of cities.
\end{problem}

\begin{code}
void solve() {
    int n, m;
    cin >> n >> m;
    vector<vector<pii>> adj(n + 1);  // (neighbour, edge id)
    rep(i, m) {
        int a, b;
        cin >> a >> b;
        adj[a].push_back({b, i});
        adj[b].push_back({a, i});
    }
    vi tin(n + 1, 0), low(n + 1, 0);
    vector<pii> bridges;
    int timer = 0;
    auto dfs = [&](auto&& self, int u, int pe) -> void {  // pe: parent edge
        tin[u] = low[u] = ++timer;
        for (auto [v, id] : adj[u]) {
            if (id == pe) continue;  // skip only the edge we came along
            if (tin[v]) low[u] = min(low[u], tin[v]);  // a back edge
            else {
                self(self, v, id);
                low[u] = min(low[u], low[v]);
                // nothing below v climbs to u or above: u-v is a bridge
                if (low[v] > tin[u]) bridges.push_back({u, v});
            }
        }
    };
    rep(v, 1, n + 1) if (!tin[v]) dfs(dfs, v, -1);
    print(sz(bridges));
    for (auto [a, b] : bridges) print(a, b);
}
\end{code}

\begin{pitfall}[Parallel edges]
For bridges, skip the edge you arrived along by its \emph{id}, not by the parent vertex. Otherwise a doubled road $u$--$v$ is
reported as a bridge, although removing one copy leaves the other. For articulation points, skipping the parent vertex is fine,
because parallel edges never change whether a vertex separates the graph.
\end{pitfall}

\begin{problem}{Necessary Cities (CSES)}
List every city whose removal would disconnect some pair of the remaining cities.
\end{problem}

\begin{code}
void solve() {
    int n, m;
    cin >> n >> m;
    vvi adj(n + 1);
    rep(i, m) {
        int a, b;
        cin >> a >> b;
        adj[a].push_back(b);
        adj[b].push_back(a);
    }
    vi tin(n + 1, 0), low(n + 1, 0);
    vector<bool> cut(n + 1, false);
    int timer = 0;
    auto dfs = [&](auto&& self, int u, int p) -> void {
        tin[u] = low[u] = ++timer;
        int children = 0;
        for (int v : adj[u]) {
            if (v == p) continue;
            if (tin[v]) low[u] = min(low[u], tin[v]);
            else {
                self(self, v, u);
                low[u] = min(low[u], low[v]);
                children++;
                // v's subtree cannot climb above u: removing u cuts it off
                if (p != 0 && low[v] >= tin[u]) cut[u] = true;
            }
        }
        if (p == 0 && children >= 2) cut[u] = true;  // the root
    };
    rep(v, 1, n + 1) if (!tin[v]) dfs(dfs, v, 0);
    vi cities;
    rep(v, 1, n + 1) if (cut[v]) cities.push_back(v);
    print(sz(cities));
    print(cities);
}
\end{code}

\usetemplate{Bridges + articulation points}{Ctrl+Alt+T, Shift+B}{cpbridges}

\subsection{Bridge trees and block-cut trees}

Removing all bridges splits a graph into \defn{2-edge-connected components}: inside one, every two vertices are joined by two
edge-disjoint paths. Contracting each of them gives the \defn{bridge tree}, whose edges are exactly the bridges. Questions like
``which edges lie on every path from $a$ to $b$?'' then become ``which edges lie on the tree path between the components of $a$ and
$b$?'', answered with LCA (Chapter~\ref{ch:trees}).

\defn{Biconnected components} (blocks) are the vertex version: maximal subgraphs without articulation points. Every edge belongs
to exactly one block, while an articulation point belongs to several. The \defn{block-cut tree} has a node for every block and
one for every articulation point, joined when the articulation point belongs to the block (Figure~\ref{fig:blockcut}). ``Which
vertices lie on every path from $a$ to $b$?'' becomes ``which articulation-point nodes lie on the tree path?''.

\begin{figure}[htbp]
\centering
\begin{tikzpicture}
  \fill[fgsky!20, rounded corners=8pt] (-0.5, -1.2) rectangle (1.65, 1.2);
  \fill[fggreen!15, rounded corners=8pt] (1.0, -1.25) rectangle (3.0, 1.25);
  \fill[fgpink!20, rounded corners=8pt] (2.25, -1.25) rectangle (4.4, -0.35);
  \node[vtx] (a) at (0, 0.8) {0}; \node[vtx] (b) at (0, -0.8) {1}; \node[vtx, hot] (c) at (1.3, 0) {2};
  \node[vtx] (d) at (2.6, 0.8) {3}; \node[vtx, hot] (e) at (2.6, -0.8) {4}; \node[vtx] (f) at (4.0, -0.8) {5};
  \draw[ed] (a) -- (b) -- (c) -- (a) (c) -- (d) -- (e) -- (c) (e) -- (f);
  \begin{scope}[xshift=6.2cm]
    \node[draw=fgblue, fill=fgsky!20, rounded corners=3pt, font=\sffamily\scriptsize, minimum height=6mm] (B1) at (0, 0)
      {$\{0, 1, 2\}$};
    \node[vtxs, hot] (C2) at (1.45, 0) {2};
    \node[draw=fggreen!70!black, fill=fggreen!15, rounded corners=3pt, font=\sffamily\scriptsize, minimum height=6mm]
      (B2) at (2.9, 0) {$\{2, 3, 4\}$};
    \node[vtxs, hot] (C4) at (4.35, 0) {4};
    \node[draw=fgpink!80!black, fill=fgpink!20, rounded corners=3pt, font=\sffamily\scriptsize, minimum height=6mm]
      (B3) at (5.6, 0) {$\{4, 5\}$};
    \draw[ed] (B1) -- (C2) -- (B2) -- (C4) -- (B3);
    \node[lbl] at (2.9, -0.85) {block-cut tree};
  \end{scope}
\end{tikzpicture}
\caption{Three blocks (shaded) that share the articulation points 2 and 4, and the block-cut tree, which alternates between
block nodes and articulation-point nodes. The single edge 4--5 is a block of its own (and a bridge).}
\label{fig:blockcut}
\end{figure}

\usetemplate{2-edge-connected components}{Ctrl+Alt+Y, B}{cp2ecc}

\usetemplate{Biconnected components (blocks)}{Ctrl+Alt+Y, Shift+B}{cpbcc}

% ======================================================================================================================
\section[Eulerian paths and circuits]{Eulerian paths and circuits \hfill \plat}\label{sec:euler}

An \defn{Eulerian path} uses every edge exactly once; an \defn{Eulerian circuit} also returns to its start. Unlike Hamiltonian
paths, which visit every \emph{vertex} once and are NP-hard to find, Eulerian paths are decided by degrees alone (plus the
condition that all edges lie in one connected component), as Table~\ref{tab:euler} shows. The reason is local: every time a walk
passes through a vertex it uses one edge in and one edge out, so all vertices except the ends must have balanced degrees.

\begin{table}[htbp]
\centering
\small
\begin{tabular}{@{}l P{4.4cm} P{6.6cm}@{}}
  \toprule
  & \thd{Circuit exists if and only if} & \thd{Path exists if and only if} \\
  \midrule
  undirected & every degree is even & exactly 0 or 2 vertices have odd degree (start at one of them) \\
  directed & in-degree = out-degree everywhere & one vertex has out $-$ in $= 1$ (the start), one has in $-$ out $= 1$ (the
    end), all others are balanced \\
  \bottomrule
\end{tabular}
\caption{Existence of Eulerian circuits and paths, assuming all edges are in one (weakly) connected component.}
\label{tab:euler}
\end{table}

\defn{Hierholzer's algorithm} constructs the tour in $\Oh{n + m}$. Walk along unused edges from the start until you get stuck; by
the degree conditions, you can only get stuck at the end vertex. Then back up along the walk: whenever a vertex on it still has
unused edges, a closed sub-tour starting there can be spliced into the walk at that point (Figure~\ref{fig:hierholzer}). With an
explicit stack this is a few lines: the top of the stack advances along an unused edge if it has one, and otherwise is popped and
appended to the output, which is the tour in reverse order.

\begin{figure}[htbp]
\centering
\begin{tikzpicture}
  \node[vtx, hot] (v1) at (0, 0) {1}; \node[vtx] (v2) at (1.6, 1.0) {2}; \node[vtx] (v3) at (3.2, 0) {3};
  \node[vtx] (v4) at (1.6, -1.0) {4}; \node[vtx] (v5) at (4.8, 1.0) {5}; \node[vtx] (v6) at (4.8, -1.0) {6};
  \draw[bed] (v1) -- node[wl] {1} (v2); \draw[bed] (v2) -- node[wl] {2} (v3); \draw[bed] (v3) -- node[wl] {3} (v4);
  \draw[bed] (v4) -- node[wl] {4} (v1);
  \draw[ged] (v3) -- node[wl] {a} (v5); \draw[ged] (v5) -- node[wl] {b} (v6); \draw[ged] (v6) -- node[wl] {c} (v3);
  \node[notel, anchor=west, text width=5.6cm] at (5.6, 0) {The first walk from 1 (blue, edges 1--4) closes without using the
    triangle at 3. Backing up to 3, which still has unused edges, starts a sub-tour a, b, c (green) that is spliced in:
    $1, 2, 3, 5, 6, 3, 4, 1$.};
\end{tikzpicture}
\caption{Hierholzer's algorithm splices closed sub-tours into the current walk wherever unused edges remain.}
\label{fig:hierholzer}
\end{figure}

\begin{problem}{Mail Delivery (CSES)}
Find a route that starts and ends at post office 1 and uses every street exactly once, or report that none exists.
\end{problem}

\begin{code}
void solve() {
    int n, m;
    cin >> n >> m;
    vector<vector<pii>> adj(n + 1);  // (neighbour, edge id)
    vi deg(n + 1, 0);
    rep(i, m) {
        int a, b;
        cin >> a >> b;
        adj[a].push_back({b, i});
        adj[b].push_back({a, i});
        deg[a]++, deg[b]++;
    }
    rep(v, 1, n + 1) if (deg[v] % 2) {  // a circuit needs even degrees
        print("IMPOSSIBLE");
        return;
    }
    // Hierholzer: walk unused edges; a vertex with none left is final
    vector<bool> used(m, false);
    vi ptr(n + 1, 0), st{1}, tour;
    while (!st.empty()) {
        int u = st.back();
        while (ptr[u] < sz(adj[u]) && used[adj[u][ptr[u]].second]) ptr[u]++;
        if (ptr[u] == sz(adj[u])) {
            tour.push_back(u);
            st.pop_back();
        } else {
            auto [v, id] = adj[u][ptr[u]];
            used[id] = true;
            st.push_back(v);
        }
    }
    if (sz(tour) != m + 1) {  // some edges are not connected to city 1
        print("IMPOSSIBLE");
        return;
    }
    print(tour);
}
\end{code}

\begin{problem}{Teleporters Path (CSES)}
Find a route from level 1 to level $n$ that uses every one-way teleporter exactly once.
\end{problem}

\begin{code}
void solve() {
    int n, m;
    cin >> n >> m;
    vvi adj(n + 1);
    vi out(n + 1, 0), in(n + 1, 0);
    rep(i, m) {
        int a, b;
        cin >> a >> b;
        adj[a].push_back(b);
        out[a]++, in[b]++;
    }
    // a path 1 -> n: out - in is +1 at 1, -1 at n, 0 elsewhere
    bool ok = out[1] - in[1] == 1 && in[n] - out[n] == 1;
    rep(v, 2, n) if (in[v] != out[v]) ok = false;
    if (!ok) {
        print("IMPOSSIBLE");
        return;
    }
    vi ptr(n + 1, 0), st{1}, path;  // Hierholzer on a directed graph
    while (!st.empty()) {
        int u = st.back();
        if (ptr[u] < sz(adj[u])) st.push_back(adj[u][ptr[u]++]);
        else path.push_back(u), st.pop_back();
    }
    if (sz(path) != m + 1) {  // some edges are unreachable from 1
        print("IMPOSSIBLE");
        return;
    }
    reverse(all(path));
    print(path);
}
\end{code}

\usetemplate{Eulerian path / circuit (Hierholzer)}{Ctrl+Alt+Y, E}{cpeuler}

\begin{trick}[Euler tours as a modelling tool]
``Use every transition exactly once'' is an Eulerian path on the graph whose edges are the transitions: de Bruijn sequences
(Section~\ref{sec:debruijn}), arranging dominoes in a chain, or reconstructing a string from all its length-$k$ substrings. And
``orient the edges so that every vertex has in-degree equal to out-degree'' is solved by following an Euler circuit after pairing
up odd vertices with dummy edges.
\end{trick}

% ======================================================================================================================
\section[Offline dynamic connectivity]{Offline dynamic connectivity \hfill \ioi}\label{sec:offline-dyncon}

Suppose edges are added \emph{and} removed, and the number of components is wanted after every operation. A DSU cannot delete,
but offline there is an elegant way around it. Every edge is alive during an interval of time. Build a segment tree over the time
axis and store each edge in the $\Oh{\log q}$ nodes that exactly cover its interval (Figure~\ref{fig:segtime}). Then walk the
segment tree depth first with a \defn{DSU with rollback} (union by size, no path compression, and a stack of the changes made):
entering a node unites its edges, reaching a leaf reads off the answer for that moment, and leaving a node undoes its unions. At
each leaf, exactly the edges alive at that time have been united. The total cost is $\Oh{q \log q \log n}$.

\begin{figure}[htbp]
\centering
\begin{tikzpicture}[x=0.9cm]
  \foreach \lv/\cnt in {0/1, 1/2, 2/4, 3/8} {
    \pgfmathsetmacro{\wdt}{8/\cnt}
    \foreach \kx in {1,...,\cnt} {
      \pgfmathsetmacro{\xl}{(\kx - 1)*\wdt}
      \draw[fgline, line width=0.5pt, fill=white] (\xl + 0.04, -\lv*0.62) rectangle (\xl + \wdt - 0.04, -\lv*0.62 - 0.46);
    }
  }
  \foreach \xl/\wdt/\lv/\cl in {1/1/3/fgorange!45, 2/2/2/fgorange!45, 4/4/1/fgorange!45, 0/2/2/fgsky!55, 5/1/3/fggreen!45,
      6/2/2/fggreen!45}
    \fill[\cl] (\xl + 0.06, -\lv*0.62 - 0.02) rectangle (\xl + \wdt - 0.06, -\lv*0.62 - 0.44);
  \foreach \tv in {0,...,7} \node[lbl] at (\tv + 0.5, -2.62) {\tv};
  \node[lbl, anchor=east] at (-0.1, -2.62) {time};
  \draw[ivl=fgorange!80] (1, -3.15) -- (8, -3.15); \node[font=\sffamily\scriptsize, anchor=west] at (8.1, -3.15) {edge $e_1$
    alive at times 1--7: three nodes};
  \draw[ivl=fgsky!80] (0, -3.55) -- (2, -3.55); \node[font=\sffamily\scriptsize, anchor=west] at (8.1, -3.55) {edge $e_2$ alive
    at times 0--1: one node};
  \draw[ivl=fggreen!70] (5, -3.95) -- (8, -3.95); \node[font=\sffamily\scriptsize, anchor=west] at (8.1, -3.95) {edge $e_3$
    alive at times 5--7: two nodes};
\end{tikzpicture}
\caption{A segment tree on time. Each edge is stored in the $\Oh{\log q}$ nodes that tile its lifetime; a leaf at time $t$ sees
exactly the edges stored on its root path, which are the edges alive at $t$.}
\label{fig:segtime}
\end{figure}

\begin{problem}{Dynamic Connectivity (CSES)}
A graph starts with $m$ edges; each of $k$ operations adds or removes an edge. Print the number of components initially and after
every operation.
\end{problem}

\begin{code}
void solve() {
    int n, m, k;
    cin >> n >> m >> k;
    // every edge is alive during an interval of times 0..k; time 0 is the
    // initial graph, time t is right after operation t
    map<pii, int> since;  // alive edge -> the time it appeared
    vector<array<int, 4>> live;  // (from, to, a, b): alive on [from, to]
    auto norm = [](int a, int b) { return pii(min(a, b), max(a, b)); };
    rep(i, m) {
        int a, b;
        cin >> a >> b;
        since[norm(a, b)] = 0;
    }
    rep(t, 1, k + 1) {
        int type, a, b;
        cin >> type >> a >> b;
        if (type == 1) since[norm(a, b)] = t;
        else {
            auto it = since.find(norm(a, b));
            live.push_back({it->second, t - 1, a, b});
            since.erase(it);
        }
    }
    for (auto [e, t] : since) live.push_back({t, k, e.first, e.second});
    // each interval is stored in O(log k) nodes of a segment tree on time
    int T = k + 1;
    vector<vector<pii>> node(4 * T);
    auto add = [&](auto&& self, int v, int l, int r, int ql, int qr, pii e) {
        if (qr < l || r < ql) return;
        if (ql <= l && r <= qr) {
            node[v].push_back(e);
            return;
        }
        int mid = (l + r) / 2;
        self(self, 2 * v, l, mid, ql, qr, e);
        self(self, 2 * v + 1, mid + 1, r, ql, qr, e);
    };
    for (auto [from, to, a, b] : live) add(add, 1, 0, T - 1, from, to, {a, b});
    // DSU with rollback: union by size, no path compression
    vi par(n + 1), siz(n + 1, 1);
    iota(all(par), 0);
    int comps = n;
    vector<pii> history;  // (attached root, new root) for undo
    auto find = [&](int x) {
        while (par[x] != x) x = par[x];
        return x;
    };
    vi ans(T);
    auto walk = [&](auto&& self, int v, int l, int r) -> void {
        int saved = sz(history);
        for (auto [a, b] : node[v]) {
            int ra = find(a), rb = find(b);
            if (ra == rb) continue;
            if (siz[ra] < siz[rb]) swap(ra, rb);
            par[rb] = ra, siz[ra] += siz[rb], comps--;
            history.push_back({rb, ra});
        }
        if (l == r) ans[l] = comps;
        else {
            int mid = (l + r) / 2;
            self(self, 2 * v, l, mid);
            self(self, 2 * v + 1, mid + 1, r);
        }
        while (sz(history) > saved) {  // undo this node's unions
            auto [rb, ra] = history.back();
            history.pop_back();
            par[rb] = rb, siz[ra] -= siz[rb], comps++;
        }
    };
    walk(walk, 1, 0, T - 1);
    print(ans);
}
\end{code}

\usetemplate{Offline dynamic connectivity (edge deletions)}{Ctrl+Alt+Y, Shift+R}{cpofflinecc}

\begin{trick}[The segment tree on time]
The technique needs only a structure that can \emph{undo} its last operation. A DSU with rollback, a knapsack array that is
copied before each insertion, a linear basis with snapshots, a stack of any kind: each item lives in $\Oh{\log q}$ nodes, and every
node's operations are undone on the way back up. It turns ``insertions and deletions, answer queries'' into ``insertions only''
for any such structure, at the cost of one logarithmic factor.
\end{trick}

% ======================================================================================================================
\section[Euler's formula and planar graphs]{Euler's formula and planar graphs \hfill \ioi}\label{sec:planar}

A graph is \emph{planar} if it can be drawn in the plane without crossing edges; such a drawing divides the plane into
\emph{faces}, including the unbounded outer face. \defn{Euler's formula} says that for a connected planar drawing,
\[
  V - E + F = 2 ,
\]
and with $C$ connected components, $V - E + F = 1 + C$ (Figure~\ref{fig:euler-formula}). Two consequences are useful. A simple
planar graph with $V \ge 3$ has $E \le 3V - 6$, so planar graphs are sparse. And counting faces, which is often what a problem
asks (``how many regions does this drawing create?''), reduces to counting vertices, edges and components, which is frequently
easier than flood filling the regions: a DSU over the vertices, applied while the edges are drawn, keeps $C$ up to date.

\begin{figure}[htbp]
\centering
\begin{tikzpicture}
  \node[vtxs] (a) at (0, 0) {}; \node[vtxs] (b) at (2, 0) {}; \node[vtxs] (c) at (4, 0) {};
  \node[vtxs] (d) at (0, 1.8) {}; \node[vtxs] (e) at (2, 1.8) {}; \node[vtxs] (f) at (4, 1.8) {};
  \draw[ed] (a) -- (b) -- (c) -- (f) -- (e) -- (d) -- (a) (b) -- (e) (a) -- (e);
  \node[font=\sffamily\small, text=fgblue] at (0.62, 1.2) {$F_1$};
  \node[font=\sffamily\small, text=fgblue] at (1.4, 0.55) {$F_2$};
  \node[font=\sffamily\small, text=fgblue] at (3, 0.9) {$F_3$};
  \node[font=\sffamily\small, text=fgblue] at (5, 0.9) {$F_4$ (outer)};
  \node[notel, anchor=west, text width=4.4cm] at (6.1, 0.9) {$V = 6$, $E = 8$, $F = 4$:\\ $6 - 8 + 4 = 2$.};
\end{tikzpicture}
\caption{Euler's formula on a small planar graph: three bounded faces and the outer face.}
\label{fig:euler-formula}
\end{figure}

% ======================================================================================================================
\section[Critical vertices and dominators]{Critical vertices, edges and dominators \hfill \ioi}\label{sec:dominators}

A vertex or edge is \emph{critical} for getting from $s$ to $t$ if every route of the relevant kind uses it. Three versions come
up, and they have completely different answers.

\begin{itemize}
  \item \textbf{Every path, undirected graph.} The critical vertices are the articulation points on the path between $s$ and $t$
        in the block-cut tree, and the critical edges are the bridges on the path in the bridge tree.
  \item \textbf{Every shortest path.} With Dijkstra from $s$ and from $t$ (on the reversed graph), an edge $(u, v, w)$ lies on
        some shortest path if and only if $d_s(u) + w + d_t(v) = D$. It lies on \emph{all} of them if moreover it carries all of
        them: $\mathit{ways}_s(u) \cdot \mathit{ways}_t(v) = \mathit{ways}_s(t)$. The counts are huge, so compare them modulo one or
        two random large primes. A vertex $v$ is critical if and only if $d_s(v) + d_t(v) = D$ and $\mathit{ways}_s(v) \cdot
        \mathit{ways}_t(v) = \mathit{ways}_s(t)$.
  \item \textbf{Every path, directed graph.} A vertex $v$ \defn{dominates} $w$ if every path from $s$ to $w$ passes through $v$.
        The immediate dominators form the \defn{dominator tree} (Figure~\ref{fig:dominator}), and the vertices critical for
        reaching $t$ are exactly the ancestors of $t$ in it. In a DAG, build it in topological order: the immediate dominator of
        $w$ is the LCA, in the tree built so far, of all predecessors of $w$. General graphs need the Lengauer--Tarjan algorithm,
        $\Oh{m\,\alpha(n)}$.
\end{itemize}

\begin{figure}[htbp]
\centering
\begin{tikzpicture}
  \begin{scope}
    \node[vtx, hot] (s) at (0, 0) {s}; \node[vtx] (a) at (1.4, 0.8) {a}; \node[vtx] (b) at (1.4, -0.8) {b};
    \node[vtx] (c) at (2.8, 0) {c}; \node[vtx] (d) at (4.2, 0) {d}; \node[vtx] (t) at (5.6, 0.8) {t};
    \node[vtx] (e) at (5.6, -0.8) {e};
    \draw[ded] (s) -- (a); \draw[ded] (s) -- (b); \draw[ded] (a) -- (c); \draw[ded] (b) -- (c); \draw[ded] (c) -- (d);
    \draw[ded] (d) -- (t); \draw[ded] (d) -- (e); \draw[ded] (e) -- (t);
    \node[lbl] at (2.8, -1.45) {a DAG from $s$};
  \end{scope}
  \begin{scope}[xshift=8.3cm, yshift=1.2cm]
    \node[vtxs, hot] (S) at (0, 0) {s};
    \node[vtxs] (A) at (-1.0, -0.8) {a}; \node[vtxs] (B) at (0, -0.8) {b}; \node[vtxs] (C) at (1.0, -0.8) {c};
    \node[vtxs] (D) at (1.0, -1.6) {d}; \node[vtxs] (T) at (0.5, -2.4) {t}; \node[vtxs] (E) at (1.5, -2.4) {e};
    \draw[ed] (S) -- (A) (S) -- (B) (S) -- (C) (C) -- (D) (D) -- (T) (D) -- (E);
    \node[lbl] at (0.2, -3.05) {dominator tree};
  \end{scope}
\end{tikzpicture}
\caption{A DAG and its dominator tree. Every path from $s$ to $t$ passes through $c$ and $d$ (the ancestors of $t$), but neither
$a$ nor $b$ is unavoidable, so $c$'s immediate dominator is $s$.}
\label{fig:dominator}
\end{figure}

\begin{recognize}
``Which roads, if closed, make the trip impossible (or longer)?'', ``checkpoints every route must pass'', ``which server failures
cut off a client''. First decide whether the question is about all paths or all \emph{shortest} paths, and whether the graph is
directed: the three cases above have different answers.
\end{recognize}

\paragraph{Link-cut trees.} Fully dynamic connectivity in a \emph{forest} (link, cut, path aggregates, online) is solved by the
\defn{link-cut tree} in amortised $\Oh{\log n}$ per operation. It represents each tree as a set of preferred paths stored in splay
trees keyed by depth, and an \emph{access} operation makes the path from the root to a vertex preferred. It is powerful but long and
error-prone; in contests, first look for an offline solution (a segment tree on time, Kruskal order, an Euler tour with a treap),
which almost always exists.

\begin{summary}[Chapter summary]
\begin{itemize}
  \item SCCs: Tarjan's single DFS with low-link values; components come out in reverse topological order of the condensation,
        ready for DP.
  \item 2-SAT: clauses become implications; unsatisfiable exactly when $x$ and $\lnot x$ share an SCC; otherwise take the literal
        whose SCC is later in topological order.
  \item Bridges: $\mathit{low}(v) > \mathit{tin}(u)$; articulation points: $\mathit{low}(v) \ge \mathit{tin}(u)$ (the root needs
        two children). Skip the parent edge by id.
  \item Eulerian paths are decided by degrees and built by Hierholzer's algorithm in linear time.
  \item Deletions offline: a segment tree on time plus any structure with undo.
  \item $V - E + F = 1 + C$ for planar drawings; dominator trees answer ``unavoidable vertex'' questions in directed graphs.
\end{itemize}
\end{summary}

\begin{practice}
\ppgroup{Strong connectivity and 2-SAT}
\begin{pplist}
  \item Flight Routes Check
  \item Planets and Kingdoms
  \item Coin Collector
  \item Giant Pizza
\end{pplist}
\ppgroup{Bridges, cut vertices and Euler tours}
\begin{pplist}
  \item Necessary Roads
  \item Necessary Cities
  \item Mail Delivery
  \item Teleporters Path
  \item De Bruijn Sequence
  \item Dynamic Connectivity
\end{pplist}
\end{practice}
